%! Minimizing Loss by Gradient Descent — Unit 8, Lesson 8.3 — Exit Ticket
\documentclass[10pt]{article}
\usepackage{linalg-article}
\usepackage{linalg-boxes}
\newcommand{\TallMath}[1]{$\displaystyle #1\rule[-1.4em]{0pt}{3.2em}$}

\begin{document}

\pageheader{Unit 8, Lesson 8.3}{Exit Ticket}
\namedateperiod

\vspace{0.15in}

No notes. Show your reasoning. Use $L(w)=(w-4)^2$ with slope $L'(w)=2w-8$.

\begin{enumerate}[leftmargin=1.8em, itemsep=16pt]
  \item \textbf{Take one step.} From $w=1$ with learning rate $s=0.25$, compute $w\leftarrow w-s\,L'(w)$.
        Give the new $w$, and check whether the loss went down by comparing $L(1)$ and $L(\text{new }w)$.
        \vspace{1.7cm}
  \item \textbf{When do we stop?} At the bottom $w=4$, what is the slope $L'(4)$, and why does that make
        gradient descent stop there? Also, in one phrase, what is the number $s$ called?
        \vspace{1.2cm}
  \item \textbf{Explain.} In one sentence, why do we \emph{subtract} $s\,L'(w)$ (step \emph{against} the
        slope) instead of adding it?
        \writelines{2}
\end{enumerate}

\end{document}
